\chapter{Question Index}
\label{chap:question_index}

\begin{tldr}
This appendix provides a comprehensive cross-referenced index of all 248 interview questions in this guide, organized by difficulty, type, frequency, and topic. Use it to target your preparation: start with the 2-hour cram session for highest-impact coverage, then expand to the 1-week plan for thorough preparation.
\end{tldr}

%==============================================================================
\section{Summary Statistics}
\label{sec:qi_summary}
%==============================================================================

\begin{table}[H]
\centering
\caption{Question distribution overview}
\begin{tabular}{lr}
\toprule
\textbf{Category} & \textbf{Count} \\
\midrule
\textbf{Total Questions} & \textbf{248} \\
\midrule
\multicolumn{2}{l}{\textit{By Difficulty}} \\
\quad L5 (Senior Engineer) & 89 \\
\quad L6 (Staff Engineer) & 127 \\
\quad L7 (Principal Engineer) & 32 \\
\midrule
\multicolumn{2}{l}{\textit{By Type}} \\
\quad Conceptual & 52 \\
\quad Trade-off & 64 \\
\quad Estimation & 28 \\
\quad Debugging & 38 \\
\quad Architecture Design & 43 \\
\quad First Principles & 23 \\
\midrule
\multicolumn{2}{l}{\textit{By Frequency}} \\
\quad \freqhigh{} (asked constantly) & 118 \\
\quad \freqmed{} (common) & 99 \\
\quad \freqlow{} (occasional) & 31 \\
\bottomrule
\end{tabular}
\end{table}

%==============================================================================
\section{Study Plans}
\label{sec:qi_study_plans}
%==============================================================================

\subsection{2-Hour Cram Session}

If you have minimal time, focus on these 20 high-frequency questions that cover the most critical topics. Each is \freqhigh{} and represents a concept that appears in the majority of ML interviews.

\begin{enumerate}
    \item Explain Flash Attention. Why is it faster even though it does the same computation? (Ch~\ref{chap:attention_transformers}, L5)
    \item Compare MHA, GQA, and MQA. When would you choose each? (Ch~\ref{chap:attention_transformers}, L5)
    \item Why does attention scale by $1/\sqrt{d_k}$? (Ch~\ref{chap:attention_transformers}, L5)
    \item What changed between GPT-2 and LLaMA architecturally? (Ch~\ref{chap:attention_transformers}, L5)
    \item Adam vs SGD with momentum---when would you pick each? (Ch~\ref{chap:training_optimization}, L5)
    \item Walk me through the RLHF pipeline for aligning an LLM. (Ch~\ref{chap:reinforcement_learning}, L5)
    \item Encoder-only vs decoder-only vs encoder-decoder---when each? (Ch~\ref{chap:nlp_architectures}, L5)
    \item Explain the bias-variance tradeoff. Does it apply to deep learning? (Ch~\ref{chap:learning_theory}, L5)
    \item When do you use softmax cross-entropy vs sigmoid BCE? (Ch~\ref{chap:loss_functions}, L5)
    \item Cross-entropy vs focal loss vs class-balanced loss for imbalanced data. (Ch~\ref{chap:loss_functions}, L5)
    \item How do you handle training-serving skew? (Ch~\ref{chap:production_systems}, L5)
    \item Your model has 95\% accuracy but stakeholders are unhappy. Why? (Ch~\ref{chap:evaluation_metrics}, L5)
    \item How do you evaluate RAG quality? What metrics matter most? (Ch~\ref{chap:long_context_rag}, L5)
    \item Long context window vs RAG---when do you choose each? (Ch~\ref{chap:long_context_rag}, L5)
    \item LoRA vs QLoRA vs full fine-tuning vs prompt tuning---decision framework. (Ch~\ref{chap:transfer_learning}, L6)
    \item Derive the gradient of softmax cross-entropy loss w.r.t.\ the logits. (Ch~\ref{chap:mathematical_foundations}, L6)
    \item Estimate the memory to train a 7B model, batch size 1, seq len 2048. (Ch~\ref{chap:attention_transformers}, L6)
    \item Data parallelism vs tensor vs pipeline parallelism---decision framework. (Ch~\ref{chap:training_optimization}, L6)
    \item PPO vs DPO vs RLHF---complete comparison for LLM alignment. (Ch~\ref{chap:reinforcement_learning}, L6)
    \item Design an LLM serving stack for 1K QPS, p99 $<$ 2s, 70B model. (Ch~\ref{chap:inference_optimization}, L7)
\end{enumerate}

\subsection{1-Week Deep Preparation}

\begin{table}[H]
\centering
\caption{1-week study plan}
\begin{tabularx}{\textwidth}{lXl}
\toprule
\textbf{Day} & \textbf{Focus Areas} & \textbf{Chapters} \\
\midrule
Day 1 & Mathematical foundations, learning theory, loss functions & Ch~\ref{chap:mathematical_foundations}--\ref{chap:loss_functions} \\
Day 2 & Attention, Transformers, embeddings & Ch~\ref{chap:attention_transformers}--\ref{chap:embeddings} \\
Day 3 & NLP architectures, vision architectures, generative models & Ch~\ref{chap:nlp_architectures}--\ref{chap:generative_models} \\
Day 4 & Retrieval, similarity, recommendations & Ch~\ref{chap:similarity_architectures}, \ref{chap:recommendation} \\
Day 5 & Training optimization, transfer learning, RLHF/alignment & Ch~\ref{chap:training_optimization}, \ref{chap:transfer_learning}, \ref{chap:reinforcement_learning} \\
Day 6 & Inference, production, evaluation, RAG, safety & Ch~\ref{chap:inference_optimization}, \ref{chap:production_systems}, \ref{chap:evaluation_metrics}, \ref{chap:long_context_rag}, \ref{chap:safety_alignment} \\
Day 7 & Mock interview: pick 20 questions across difficulty levels & All chapters \\
\bottomrule
\end{tabularx}
\end{table}

\textbf{Day 7 mock interview strategy}: Select 8 L5 questions, 8 L6 questions, and 4 L7 questions. Time yourself: 3--5 minutes per L5, 5--8 minutes per L6, 8--10 minutes per L7. Record yourself and review for clarity, structure, and technical depth.

\subsection{By Company Focus Area}

\begin{table}[H]
\centering
\caption{Company focus area to chapter mapping}
\begin{tabularx}{\textwidth}{lX}
\toprule
\textbf{Team/Role Focus} & \textbf{Priority Chapters} \\
\midrule
LLM / Foundation Models & Ch~\ref{chap:attention_transformers}, \ref{chap:nlp_architectures}, \ref{chap:efficient_architectures}, \ref{chap:training_optimization}, \ref{chap:reinforcement_learning}, \ref{chap:inference_optimization}, \ref{chap:safety_alignment} \\
Ads / Ranking / RecSys & Ch~\ref{chap:loss_functions}, \ref{chap:embeddings}, \ref{chap:similarity_architectures}, \ref{chap:recommendation}, \ref{chap:production_systems}, \ref{chap:evaluation_metrics} \\
Vision / Multimodal & Ch~\ref{chap:vision_architectures}, \ref{chap:generative_models}, \ref{chap:multimodal_learning}, \ref{chap:self_supervised} \\
ML Infrastructure & Ch~\ref{chap:efficient_architectures}, \ref{chap:training_optimization}, \ref{chap:production_systems}, \ref{chap:inference_optimization} \\
Search / Retrieval & Ch~\ref{chap:similarity_architectures}, \ref{chap:embeddings}, \ref{chap:long_context_rag}, \ref{chap:evaluation_metrics} \\
Safety / Alignment & Ch~\ref{chap:reinforcement_learning}, \ref{chap:safety_alignment}, \ref{chap:long_context_rag} \\
\bottomrule
\end{tabularx}
\end{table}

%==============================================================================
\section{Questions by Difficulty}
\label{sec:qi_by_difficulty}
%==============================================================================

\subsection{L5 Questions (Senior Engineer)}

89 questions. These test solid fundamentals, clear communication, and the ability to compare approaches.

\begin{longtable}{rp{8cm}lr}
\caption{L5 (Senior) questions} \\
\toprule
\textbf{\#} & \textbf{Question} & \textbf{Ch} & \textbf{Type} \\
\midrule
\endfirsthead
\toprule
\textbf{\#} & \textbf{Question} & \textbf{Ch} & \textbf{Type} \\
\midrule
\endhead
\midrule
\multicolumn{4}{r}{\textit{Continued on next page}} \\
\endfoot
\bottomrule
\endlastfoot

% Ch01 - Mathematical Foundations
1 & Explain SVD and give three applications in ML. \freqhigh & \ref{chap:mathematical_foundations} & Conceptual \\
2 & What is KL divergence? Why is it not a true distance metric? \freqhigh & \ref{chap:mathematical_foundations} & Conceptual \\
3 & Estimate memory for a 50K$\times$768 embedding matrix (FP32/FP16/INT8). \freqhigh & \ref{chap:mathematical_foundations} & Estimation \\
4 & What are eigenvalues and why do they matter for training? \freqmed & \ref{chap:mathematical_foundations} & Conceptual \\
5 & Why does the chain rule matter for backprop? Walk through 3 layers. \freqhigh & \ref{chap:mathematical_foundations} & First Principles \\
6 & Estimate the FLOPs for multiplying (1024,4096)$\times$(4096,1024). \freqmed & \ref{chap:mathematical_foundations} & Estimation \\

% Ch02 - Learning Theory
7 & Explain double descent. Why does test error decrease in overparameterized regime? \freqhigh & \ref{chap:learning_theory} & Conceptual \\
8 & Explain the bias-variance tradeoff. Does it apply to deep learning? \freqhigh & \ref{chap:learning_theory} & Conceptual \\

% Ch03 - Loss Functions
9 & When do you use softmax cross-entropy vs sigmoid BCE? \freqhigh & \ref{chap:loss_functions} & Conceptual \\
10 & Why does SimCLR need large batch sizes, and alternatives? \freqhigh & \ref{chap:loss_functions} & Conceptual \\
11 & 95\% accuracy but loss still decreasing. Should you keep training? \freqhigh & \ref{chap:loss_functions} & Trade-off \\
12 & Product recommendation with implicit feedback---what loss? \freqhigh & \ref{chap:loss_functions} & Arch Design \\
13 & CE vs focal loss vs class-balanced for imbalanced classification. \freqhigh & \ref{chap:loss_functions} & Trade-off \\
14 & Embedding model for semantic search across 100M docs---which loss? \freqhigh & \ref{chap:loss_functions} & Arch Design \\

% Ch04 - Activation & Normalization
15 & Training loss is NaN after a few hundred steps. Diagnose. \freqhigh & \ref{chap:activation_normalization} & Debugging \\
16 & Why GELU over ReLU? Why SwiGLU over GELU? \freqhigh & \ref{chap:activation_normalization} & Trade-off \\
17 & RMSNorm vs LayerNorm---when would you choose each? \freqmed & \ref{chap:activation_normalization} & Trade-off \\
18 & Explain the dying ReLU problem. How do modern activations solve it? \freqhigh & \ref{chap:activation_normalization} & Conceptual \\
19 & Validation accuracy lower than training---normalization debugging. \freqmed & \ref{chap:activation_normalization} & Debugging \\

% Ch05 - Attention & Transformers
20 & Explain Flash Attention. Why is it faster with the same computation? \freqhigh & \ref{chap:attention_transformers} & Conceptual \\
21 & How does RoPE encode position information? Why preferred? \freqhigh & \ref{chap:attention_transformers} & Conceptual \\
22 & Compare MHA, GQA, and MQA. When would you choose each? \freqhigh & \ref{chap:attention_transformers} & Trade-off \\
23 & Why does attention scale by $1/\sqrt{d_k}$? What without it? \freqhigh & \ref{chap:attention_transformers} & First Principles \\
24 & Why is self-attention permutation equivariant? Why does it matter? \freqmed & \ref{chap:attention_transformers} & First Principles \\
25 & Pre-norm vs Post-norm Transformers---which and why? \freqmed & \ref{chap:attention_transformers} & Trade-off \\
26 & What changed between GPT-2 (2019) and LLaMA (2023)? Why each change? \freqhigh & \ref{chap:attention_transformers} & Conceptual \\

% Ch06 - Embeddings
27 & Static vs contextual embeddings---when still use static? \freqmed & \ref{chap:embeddings} & Trade-off \\
28 & Estimate memory for 50K vocab, 768-dim embedding table (FP16/INT8). \freqmed & \ref{chap:embeddings} & Estimation \\
29 & How does BPE interact with embedding quality? \freqhigh & \ref{chap:embeddings} & Conceptual \\
30 & How would you debug poor embedding quality? \freqhigh & \ref{chap:embeddings} & Debugging \\
31 & Build a semantic search system---how choose and evaluate embeddings? \freqhigh & \ref{chap:embeddings} & Arch Design \\

% Ch07 - Similarity & Metric Learning
32 & Two-tower retrieves keyword matches but not relevant docs. Fix? \freqhigh & \ref{chap:similarity_architectures} & Debugging \\
33 & Siamese vs Triplet vs Two-Tower---decision framework. \freqhigh & \ref{chap:similarity_architectures} & Trade-off \\
34 & Cross-encoder vs bi-encoder---when is quality gap worth latency? \freqhigh & \ref{chap:similarity_architectures} & Trade-off \\
35 & Cosine similarity vs dot product vs L2 distance---when each? \freqhigh & \ref{chap:similarity_architectures} & Conceptual \\

% Ch08 - Self-Supervised Learning
36 & What makes a good augmentation policy for contrastive learning? \freqmed & \ref{chap:self_supervised} & Conceptual \\
37 & SimCLR vs BYOL vs MAE---decision framework. \freqhigh & \ref{chap:self_supervised} & Trade-off \\
38 & BERT vs GPT pre-training objectives---fundamental tradeoffs. \freqhigh & \ref{chap:self_supervised} & Trade-off \\
39 & Why do non-contrastive methods (BYOL) work without negatives? \freqhigh & \ref{chap:self_supervised} & Conceptual \\

% Ch09 - NLP Architectures
40 & Encoder-only vs decoder-only vs encoder-decoder---when each? \freqhigh & \ref{chap:nlp_architectures} & Trade-off \\
41 & BPE vs SentencePiece vs character-level---tradeoffs. \freqmed & \ref{chap:nlp_architectures} & Trade-off \\

% Ch10 - Vision Architectures
42 & CNN vs ViT for 10K labeled images? What changes with 10M? \freqhigh & \ref{chap:vision_architectures} & Trade-off \\
43 & What is the effective receptive field and why does it matter? \freqmed & \ref{chap:vision_architectures} & Conceptual \\
44 & How DETR works and why it is a paradigm shift in detection. \freqhigh & \ref{chap:vision_architectures} & Conceptual \\

% Ch11 - Generative Models
45 & Reparameterization trick: why needed for VAE training? \freqhigh & \ref{chap:generative_models} & First Principles \\
46 & What is classifier-free guidance and why it improves diffusion quality? \freqhigh & \ref{chap:generative_models} & Conceptual \\
47 & ELBO derivation for VAEs and why blurry images. Fixes? \freqhigh & \ref{chap:generative_models} & First Principles \\
48 & Stable Diffusion end-to-end: text prompt to generated image. \freqhigh & \ref{chap:generative_models} & Conceptual \\

% Ch12 - Recommendation Systems
49 & How do you handle the cold start problem? \freqhigh & \ref{chap:recommendation} & Conceptual \\
50 & Two-Tower vs interaction-based for candidate generation. \freqhigh & \ref{chap:recommendation} & Trade-off \\
51 & Why do recommendation models need calibration? \freqhigh & \ref{chap:recommendation} & Conceptual \\

% Ch13 - Graph Neural Networks
52 & GNN performance degrades with more layers. What and how to fix? \freqmed & \ref{chap:graph_neural_networks} & Debugging \\
53 & GCN vs GraphSAGE vs GAT---when would you use each? \freqmed & \ref{chap:graph_neural_networks} & Trade-off \\
54 & How do you handle heterogeneous graphs in a GNN? \freqmed & \ref{chap:graph_neural_networks} & Conceptual \\

% Ch15 - Training Optimization
55 & Model's loss plateaus after initial descent. Debugging process. \freqhigh & \ref{chap:training_optimization} & Debugging \\
56 & Adam vs SGD with momentum---when would you pick each? \freqhigh & \ref{chap:training_optimization} & Trade-off \\
57 & Explain gradient accumulation and when to use it. \freqmed & \ref{chap:training_optimization} & Conceptual \\
58 & Training loss is NaN after 1000 steps. Debugging process. \freqhigh & \ref{chap:training_optimization} & Debugging \\
59 & Cosine decay vs constant LR with warmup---when each? \freqmed & \ref{chap:training_optimization} & Trade-off \\

% Ch16 - Transfer Learning
60 & Estimate trainable params for LoRA rank 16 on a 7B model. \freqhigh & \ref{chap:transfer_learning} & Estimation \\
61 & What is negative transfer? How to detect and prevent it? \freqmed & \ref{chap:transfer_learning} & Conceptual \\
62 & Adapter layers vs LoRA vs prefix tuning---differences and when each? \freqmed & \ref{chap:transfer_learning} & Trade-off \\
63 & Fine-tune 70B LLM on single A100 80GB---what approach? \freqhigh & \ref{chap:transfer_learning} & Arch Design \\
64 & Compare LoRA, full fine-tuning, and prompt tuning. When each? \freqhigh & \ref{chap:transfer_learning} & Trade-off \\

% Ch17 - Production Systems
65 & How do you handle training-serving skew? \freqhigh & \ref{chap:production_systems} & Conceptual \\
66 & Shadow deployment vs A/B test vs multi-armed bandit---when each? \freqhigh & \ref{chap:production_systems} & Trade-off \\
67 & How do you detect concept drift vs data drift? \freqmed & \ref{chap:production_systems} & Conceptual \\

% Ch18 - Evaluation Metrics
68 & 95\% accuracy but stakeholders unhappy. What is going on? \freqhigh & \ref{chap:evaluation_metrics} & Conceptual \\
69 & How do you evaluate a search ranking system? \freqhigh & \ref{chap:evaluation_metrics} & Conceptual \\
70 & Precision 0.95 but recall 0.30. What do you do? \freqhigh & \ref{chap:evaluation_metrics} & Trade-off \\
71 & How do you know if improvement is statistically significant? \freqhigh & \ref{chap:evaluation_metrics} & Conceptual \\
72 & NDCG vs MAP vs MRR for ranking---when each? \freqhigh & \ref{chap:evaluation_metrics} & Trade-off \\
73 & What is wrong with accuracy as a metric? When is it fine? \freqhigh & \ref{chap:evaluation_metrics} & Conceptual \\

% Ch19 - Decision Frameworks
74 & GPU vs TPU vs CPU for inference---decision framework. \freqmed & \ref{chap:decision_frameworks} & Trade-off \\
75 & Would you use BERT or GPT for sentiment analysis? \freqhigh & \ref{chap:decision_frameworks} & Trade-off \\
76 & Gradient boosting vs neural networks for tabular data? \freqhigh & \ref{chap:decision_frameworks} & Trade-off \\

% Ch20 - Reinforcement Learning
77 & Walk me through the RLHF pipeline for aligning an LLM. \freqhigh & \ref{chap:reinforcement_learning} & Conceptual \\
78 & What is reward hacking? 3 examples and mitigation. \freqhigh & \ref{chap:reinforcement_learning} & Conceptual \\
79 & Why does the KL penalty matter in RLHF? Without it? \freqhigh & \ref{chap:reinforcement_learning} & First Principles \\

% Ch21 - Multimodal Learning
80 & How does CLIP's contrastive loss create a shared embedding space? \freqmed & \ref{chap:multimodal_learning} & First Principles \\
81 & CLIP vs SigLIP---what changed and why? \freqmed & \ref{chap:multimodal_learning} & Conceptual \\
82 & Compare early, late, and cross-attention fusion. When each? \freqhigh & \ref{chap:multimodal_learning} & Conceptual \\

% Ch22 - Inference Optimization
83 & Explain speculative decoding. When use it and when not? \freqhigh & \ref{chap:inference_optimization} & Conceptual \\

% Ch23 - Safety & Alignment
84 & Jailbreak attacks: how do they work and how defend? \freqhigh & \ref{chap:safety_alignment} & Conceptual \\
85 & How do you detect and reduce hallucination in production LLM? \freqhigh & \ref{chap:safety_alignment} & Conceptual \\

% Ch24 - Long Context & RAG
86 & How do you evaluate RAG quality? What metrics matter most? \freqhigh & \ref{chap:long_context_rag} & Conceptual \\
87 & Chunking strategy: fixed-size vs semantic vs document-aware. \freqhigh & \ref{chap:long_context_rag} & Trade-off \\
88 & Long context window vs RAG---when do you choose each? \freqhigh & \ref{chap:long_context_rag} & Trade-off \\
89 & Speculative decoding: when use and when not? \freqhigh & \ref{chap:inference_optimization} & Conceptual \\
\end{longtable}

\subsection{L6 Questions (Staff Engineer)}

127 questions. These test depth, tradeoff reasoning, estimation skills, and debugging methodology.

\begin{longtable}{rp{8cm}lr}
\caption{L6 (Staff) questions} \\
\toprule
\textbf{\#} & \textbf{Question} & \textbf{Ch} & \textbf{Type} \\
\midrule
\endfirsthead
\toprule
\textbf{\#} & \textbf{Question} & \textbf{Ch} & \textbf{Type} \\
\midrule
\endhead
\midrule
\multicolumn{4}{r}{\textit{Continued on next page}} \\
\endfoot
\bottomrule
\endlastfoot

% Ch01
1 & Derive the gradient of softmax CE loss w.r.t.\ logits. \freqhigh & \ref{chap:mathematical_foundations} & First Princ. \\
2 & KL divergence asymmetry---when does it matter in practice? \freqmed & \ref{chap:mathematical_foundations} & Conceptual \\
3 & Gradient norms 1000$\times$ larger in early layers. What is happening? \freqmed & \ref{chap:mathematical_foundations} & Debugging \\

% Ch02
4 & Fixed compute budget---model size vs training data decision. \freqhigh & \ref{chap:learning_theory} & Trade-off \\
5 & Are emergent abilities in LLMs real or artifact? \freqmed & \ref{chap:learning_theory} & Conceptual \\
6 & 10B model overparameterized for your task. Should you reduce it? \freqhigh & \ref{chap:learning_theory} & Trade-off \\
7 & Estimate compute-optimal model for $10^{22}$ FLOPs (Chinchilla). \freqhigh & \ref{chap:learning_theory} & Estimation \\
8 & What is grokking and what does it tell us about generalization? \freqmed & \ref{chap:learning_theory} & Conceptual \\
9 & Training loss decreasing, val loss plateaued, 10$\times$ unlabeled data. \freqmed & \ref{chap:learning_theory} & Trade-off \\

% Ch03
10 & Visual search for 10M products---what loss function? \freqhigh & \ref{chap:loss_functions} & Arch Design \\
11 & Validation miscalibrated despite decreasing train loss. Why? \freqmed & \ref{chap:loss_functions} & Debugging \\
12 & Contrastive loss jumps mid-training and doesn't recover. Diagnose. \freqlow & \ref{chap:loss_functions} & Debugging \\
13 & Pointwise vs pairwise vs listwise ranking losses---when each? \freqhigh & \ref{chap:loss_functions} & Trade-off \\
14 & Retrieval with click data (implicit feedback)---design the loss. \freqhigh & \ref{chap:loss_functions} & Arch Design \\

% Ch04
15 & Why don't Transformers use BatchNorm? First principles. \freqmed & \ref{chap:activation_normalization} & First Princ. \\
16 & Different behavior at train vs inference---normalization issue? \freqmed & \ref{chap:activation_normalization} & Debugging \\
17 & Derive why BatchNorm doesn't work well for Transformers. \freqmed & \ref{chap:activation_normalization} & First Princ. \\

% Ch05
18 & Estimate memory to train 7B model (BS=1, seq=2048). \freqhigh & \ref{chap:attention_transformers} & Estimation \\
19 & Calculate FLOPs for single forward pass through 7B Transformer. \freqhigh & \ref{chap:attention_transformers} & Estimation \\
20 & Calculate KV cache memory for LLaMA 70B at 128K context. \freqhigh & \ref{chap:attention_transformers} & Estimation \\
21 & Degrading perplexity at 32K despite 8K training. Diagnose/fix. \freqmed & \ref{chap:attention_transformers} & Debugging \\
22 & Compare Transformer vs SSM (Mamba)---when choose each? \freqmed & \ref{chap:attention_transformers} & Trade-off \\
23 & Flash Attention vs Linear Attention---when each? \freqmed & \ref{chap:attention_transformers} & Trade-off \\

% Ch06
24 & Embedding space: unrelated items cluster together. Diagnose/fix. \freqmed & \ref{chap:embeddings} & Debugging \\
25 & Design embedding system for 100M product catalog. \freqhigh & \ref{chap:embeddings} & Arch Design \\
26 & What is embedding collapse? How detect/prevent it? \freqmed & \ref{chap:embeddings} & Conceptual \\

% Ch07
27 & Design multi-stage retrieval pipeline (generation/rerank/score). \freqhigh & \ref{chap:similarity_architectures} & Arch Design \\
28 & Estimate ANN index size and query latency for 1B 768-dim embeddings. \freqmed & \ref{chap:similarity_architectures} & Estimation \\
29 & Two-tower recall@10 good but precision@10 poor. What's wrong? \freqmed & \ref{chap:similarity_architectures} & Debugging \\
30 & Hard negative mining strategies---choose and failure modes. \freqmed & \ref{chap:similarity_architectures} & Trade-off \\
31 & Estimate QPS for HNSW on 100M documents. \freqmed & \ref{chap:similarity_architectures} & Estimation \\
32 & Estimate QPS: HNSW on 100M + cross-encoder reranker on top 100. \freqmed & \ref{chap:similarity_architectures} & Estimation \\

% Ch08
33 & Design SSL pre-training pipeline for 10M unlabeled images. \freqhigh & \ref{chap:self_supervised} & Arch Design \\
34 & SSL downstream plateaus despite 10$\times$ more pre-training data. Why? \freqmed & \ref{chap:self_supervised} & Debugging \\
35 & Estimate compute cost of pre-training BERT-base on 16GB text. \freqmed & \ref{chap:self_supervised} & Estimation \\
36 & Next-token vs MLM---why did autoregressive win for generative AI? \freqhigh & \ref{chap:self_supervised} & First Princ. \\
37 & Medical imaging: 10M unlabeled, 10K labeled---SSL pipeline. \freqhigh & \ref{chap:self_supervised} & Arch Design \\
38 & InfoNCE loss and mutual information. Why temperature matters? \freqmed & \ref{chap:self_supervised} & First Princ. \\

% Ch09
39 & Fine-tuned LLM fluent but factually wrong. Diagnose/fix. \freqhigh & \ref{chap:nlp_architectures} & Debugging \\
40 & Estimate inference latency: 100 tokens, 7B model, single A100. \freqmed & \ref{chap:nlp_architectures} & Estimation \\
41 & Why decoder-only dominates despite encoder-decoder flexibility? \freqhigh & \ref{chap:nlp_architectures} & Conceptual \\
42 & Walk through DPO loss---how it eliminates separate reward model. \freqhigh & \ref{chap:nlp_architectures} & First Princ. \\
43 & Customer support ticket classification (50 categories, multilingual). \freqhigh & \ref{chap:nlp_architectures} & Arch Design \\
44 & Compare RoPE vs learned vs sinusoidal positional embeddings. \freqmed & \ref{chap:nlp_architectures} & Conceptual \\

% Ch10
45 & Estimate FLOPs: ResNet-50 vs ViT-B/16 on 224$\times$224. \freqmed & \ref{chap:vision_architectures} & Estimation \\
46 & Edge device vision backbone ($<$10W power budget). \freqmed & \ref{chap:vision_architectures} & Arch Design \\
47 & ViT 92\% on ImageNet but 65\% on your domain. Diagnose. \freqmed & \ref{chap:vision_architectures} & Debugging \\
48 & Visual search: 100M images, sub-200ms latency. Architecture. \freqmed & \ref{chap:vision_architectures} & Arch Design \\
49 & ResNet skip vs DenseNet dense vs Swin shifted windows. \freqmed & \ref{chap:vision_architectures} & Trade-off \\

% Ch11
50 & Diffusion vs GAN vs VAE vs autoregressive---comparison/framework. \freqhigh & \ref{chap:generative_models} & Trade-off \\
51 & Diffusion mode collapse for certain categories. Diagnose/fix. \freqmed & \ref{chap:generative_models} & Debugging \\
52 & Product photography pipeline---image generation design. \freqmed & \ref{chap:generative_models} & Arch Design \\
53 & WGAN loss: Wasserstein distance and GAN stability. \freqmed & \ref{chap:generative_models} & First Princ. \\

% Ch12
54 & Design rec system for video streaming (100M users, 1M videos). \freqhigh & \ref{chap:recommendation} & Arch Design \\
55 & Great offline metrics but poor online performance. What went wrong? \freqhigh & \ref{chap:recommendation} & Debugging \\
56 & Design ranking model for e-commerce product ads. \freqhigh & \ref{chap:recommendation} & Arch Design \\
57 & Estimate memory for DLRM embeddings (100M users, 10M items). \freqhigh & \ref{chap:recommendation} & Estimation \\
58 & Ads ranking: score 100 candidates in $<$50ms---constraints? \freqmed & \ref{chap:recommendation} & Estimation \\
59 & CTR AUC +0.5\% but online revenue $-$3\%. Diagnose. \freqhigh & \ref{chap:recommendation} & Debugging \\
60 & Shared-bottom vs MMoE vs PLE for multi-task ads. \freqmed & \ref{chap:recommendation} & Trade-off \\
61 & ESMM: sample selection bias in CVR prediction. \freqmed & \ref{chap:recommendation} & Conceptual \\

% Ch13
62 & GNN-based fraud detection on financial transaction graph. \freqmed & \ref{chap:graph_neural_networks} & Arch Design \\
63 & When should you NOT use a GNN? Simpler alternatives? \freqmed & \ref{chap:graph_neural_networks} & Trade-off \\
64 & GCN vs GraphSAGE vs GAT for production rec system at scale? \freqmed & \ref{chap:graph_neural_networks} & Trade-off \\
65 & Estimate memory for 3-layer GCN on 10M-node graph. \freqlow & \ref{chap:graph_neural_networks} & Estimation \\
66 & WL expressivity limitations of message-passing GNNs. \freqlow & \ref{chap:graph_neural_networks} & First Princ. \\

% Ch14
67 & On-device ML architecture (phone, $<$100ms, $<$50MB). \freqhigh & \ref{chap:efficient_architectures} & Arch Design \\
68 & MoE vs dense scaling---when choose each? \freqmed & \ref{chap:efficient_architectures} & Trade-off \\
69 & 100K-token documents: Transformer vs Mamba. \freqmed & \ref{chap:efficient_architectures} & Trade-off \\
70 & Estimate MoE compute savings (8 experts, top-2 vs dense). \freqmed & \ref{chap:efficient_architectures} & Estimation \\
71 & INT8 vs INT4 vs FP8 quantization---decision framework. \freqhigh & \ref{chap:efficient_architectures} & Trade-off \\
72 & Quantized model fails on code/math but benchmarks fine. Why? \freqmed & \ref{chap:efficient_architectures} & Debugging \\
73 & KD pipeline: compress 70B teacher to 7B student. \freqmed & \ref{chap:efficient_architectures} & Arch Design \\
74 & MoE routing collapse: detect and fix. \freqmed & \ref{chap:efficient_architectures} & Debugging \\
75 & Structured vs unstructured pruning---which gives real speedup? \freqmed & \ref{chap:efficient_architectures} & Trade-off \\

% Ch15
76 & Scale training of 70B model across 256 GPUs. \freqhigh & \ref{chap:training_optimization} & Arch Design \\
77 & Loss spikes every N steps then recovers. Diagnose. \freqmed & \ref{chap:training_optimization} & Debugging \\
78 & Distributed training 60\% as fast as expected with 8 GPUs. Why? \freqmed & \ref{chap:training_optimization} & Debugging \\
79 & GPU memory for 13B model: LoRA vs full fine-tuning. \freqhigh & \ref{chap:training_optimization} & Estimation \\
80 & Data parallelism vs tensor vs pipeline parallelism---framework. \freqhigh & \ref{chap:training_optimization} & Trade-off \\
81 & FP16 vs BF16 vs FP8 for training---when each? \freqmed & \ref{chap:training_optimization} & Trade-off \\

% Ch16
82 & LoRA fine-tuned model catastrophically forgets general capabilities. \freqhigh & \ref{chap:transfer_learning} & Debugging \\
83 & LoRA vs QLoRA vs full FT vs prompt tuning---decision framework. \freqhigh & \ref{chap:transfer_learning} & Trade-off \\
84 & Design fine-tuning pipeline for foundation model, 10K examples. \freqmed & \ref{chap:transfer_learning} & Arch Design \\
85 & Why does LoRA work? Low-rank assumption about fine-tuning. \freqmed & \ref{chap:transfer_learning} & First Princ. \\
86 & Fine-tuned model: good on test set, poor in production. Why? \freqhigh & \ref{chap:transfer_learning} & Debugging \\

% Ch17
87 & Design ML infra for search ranking at 10K QPS. \freqhigh & \ref{chap:production_systems} & Arch Design \\
88 & CTR dropped 5\% overnight. Walk through diagnosis. \freqhigh & \ref{chap:production_systems} & Debugging \\
89 & Set up A/B test for new recommendation model. \freqhigh & \ref{chap:production_systems} & Arch Design \\
90 & Online metrics diverge from offline after 2 weeks. Why? \freqhigh & \ref{chap:production_systems} & Debugging \\
91 & Design monitoring system for ML model in production. \freqmed & \ref{chap:production_systems} & Arch Design \\
92 & Feature store: batch vs streaming features---architecture. \freqmed & \ref{chap:production_systems} & Arch Design \\

% Ch18
93 & AUC 0.95 but ECE 0.15 (poor calibration). Why and how fix? \freqhigh & \ref{chap:evaluation_metrics} & Debugging \\
94 & Estimate sample size for A/B test detecting 1\% CTR improvement. \freqmed & \ref{chap:evaluation_metrics} & Estimation \\
95 & Offline metrics up but online metrics down. 5 possible reasons. \freqhigh & \ref{chap:evaluation_metrics} & Debugging \\

% Ch19
96 & Design content recommendation for a news app. \freqhigh & \ref{chap:decision_frameworks} & Arch Design \\
97 & Design ML architecture for content moderation system. \freqhigh & \ref{chap:decision_frameworks} & Arch Design \\
98 & Latency 50ms, need $<$10ms. Optimization decision tree. \freqhigh & \ref{chap:decision_frameworks} & Trade-off \\
99 & 3 months, 3 engineers: fine-tune LLM or train custom model? \freqhigh & \ref{chap:decision_frameworks} & Trade-off \\

% Ch20
100 & PPO vs DPO vs RLHF---complete comparison for alignment. \freqhigh & \ref{chap:reinforcement_learning} & Trade-off \\
101 & RLHF-tuned model becomes sycophantic. Diagnose and fix. \freqhigh & \ref{chap:reinforcement_learning} & Debugging \\
102 & Estimate memory for RLHF (4 models: policy/ref/RM/critic). \freqmed & \ref{chap:reinforcement_learning} & Estimation \\
103 & Design reward model training pipeline. What data needed? \freqmed & \ref{chap:reinforcement_learning} & Arch Design \\
104 & Constitutional AI vs RLHF vs DPO---when each? \freqmed & \ref{chap:reinforcement_learning} & Trade-off \\
105 & Process vs outcome reward models---tradeoffs. \freqmed & \ref{chap:reinforcement_learning} & Trade-off \\

% Ch21
106 & Design multimodal content understanding for social media. \freqhigh & \ref{chap:multimodal_learning} & Arch Design \\
107 & Adapter-based (LLaVA) vs natively multimodal (Gemini). \freqhigh & \ref{chap:multimodal_learning} & Trade-off \\
108 & VLM hallucinates visual details not in the image. Why/fix? \freqhigh & \ref{chap:multimodal_learning} & Debugging \\
109 & Visual tokens use 576 tokens of context. Reduce cost? \freqmed & \ref{chap:multimodal_learning} & Trade-off \\
110 & Multimodal product search (text or image queries). \freqmed & \ref{chap:multimodal_learning} & Arch Design \\
111 & Add audio understanding to existing VLM. \freqlow & \ref{chap:multimodal_learning} & Arch Design \\

% Ch22
112 & KV cache memory bottleneck---dropping requests. What to do? \freqhigh & \ref{chap:inference_optimization} & Debugging \\
113 & INT4 memory savings for 70B model. Quality tradeoff? \freqhigh & \ref{chap:inference_optimization} & Estimation \\
114 & GPTQ vs AWQ vs GGUF quantization---when each? \freqmed & \ref{chap:inference_optimization} & Trade-off \\
115 & Tensor parallelism vs pipeline parallelism for inference. \freqhigh & \ref{chap:inference_optimization} & Trade-off \\
116 & p99 latency 10$\times$ the p50. Diagnose. \freqmed & \ref{chap:inference_optimization} & Debugging \\
117 & INT4 quantized model: garbage on some prompts, fine on others. \freqmed & \ref{chap:inference_optimization} & Debugging \\
118 & Speculative decoding slower than standard. What went wrong? \freqlow & \ref{chap:inference_optimization} & Debugging \\

% Ch23
119 & Content filtering pipeline for LLM chatbot, $<$50ms latency. \freqhigh & \ref{chap:safety_alignment} & Arch Design \\
120 & What is mechanistic interpretability? Concrete circuit example. \freqmed & \ref{chap:safety_alignment} & Conceptual \\
121 & Hallucination detection system design. \freqmed & \ref{chap:safety_alignment} & Arch Design \\
122 & Red teaming program for LLM product. \freqlow & \ref{chap:safety_alignment} & Arch Design \\

% Ch24
123 & Estimate cost per query: RAG vs long-context LLM (128K). \freqhigh & \ref{chap:long_context_rag} & Estimation \\
124 & RAG returns relevant docs but LLM ignores them. Why? \freqhigh & \ref{chap:long_context_rag} & Debugging \\
125 & HyDE vs query rewriting vs query decomposition---when each? \freqmed & \ref{chap:long_context_rag} & Trade-off \\
126 & Design RAG for enterprise doc Q\&A over 10M documents. \freqhigh & \ref{chap:long_context_rag} & Arch Design \\
127 & Cross-encoder vs ColBERT vs bi-encoder reranking. \freqmed & \ref{chap:long_context_rag} & Trade-off \\
\end{longtable}

\subsection{L7 Questions (Principal Engineer)}

32 questions. These test architectural vision, systems thinking, novel problem decomposition, and second-order effects.

\begin{longtable}{rp{8cm}lr}
\caption{L7 (Principal) questions} \\
\toprule
\textbf{\#} & \textbf{Question} & \textbf{Ch} & \textbf{Type} \\
\midrule
\endfirsthead
\toprule
\textbf{\#} & \textbf{Question} & \textbf{Ch} & \textbf{Type} \\
\midrule
\endhead
\bottomrule
\endlastfoot

1 & Rank of attention matrix and efficiency implications. \freqlow & \ref{chap:mathematical_foundations} & First Princ. \\
2 & Neural scaling laws: what scales, what doesn't, when they break. \freqmed & \ref{chap:learning_theory} & Conceptual \\
3 & Why do LLMs generalize with more params than training examples? \freqmed & \ref{chap:learning_theory} & First Princ. \\
4 & Multi-task model (classification + regression + ranking)---loss design. \freqmed & \ref{chap:loss_functions} & Arch Design \\
5 & Design attention mechanism for 1M-token documents. \freqlow & \ref{chap:attention_transformers} & Arch Design \\
6 & Attention patterns: all heads attend to same positions. Why? \freqlow & \ref{chap:attention_transformers} & Debugging \\
7 & Attention FLOPs vs FFN FLOPs at different sequence lengths. \freqlow & \ref{chap:attention_transformers} & Estimation \\
8 & Modify Transformer for real-time streaming (live transcription). \freqlow & \ref{chap:attention_transformers} & Arch Design \\
9 & Contrastive model: good for frequent items, poor for rare. Why? \freqlow & \ref{chap:similarity_architectures} & Debugging \\
10 & Update embeddings in production without full ANN re-index. \freqlow & \ref{chap:similarity_architectures} & Arch Design \\
11 & Multilingual doc understanding (50+ languages, mixed, tables). \freqmed & \ref{chap:nlp_architectures} & Arch Design \\
12 & How does in-context learning work mechanistically? \freqlow & \ref{chap:nlp_architectures} & First Princ. \\
13 & Estimate inference cost for 512$\times$512 Stable Diffusion on A100. \freqlow & \ref{chap:generative_models} & Estimation \\
14 & Flow matching vs denoising diffusion---what changed and why? \freqlow & \ref{chap:generative_models} & Conceptual \\
15 & Design real-time personalized notification system. \freqmed & \ref{chap:recommendation} & Arch Design \\
16 & Embedding infrastructure for rec system with 1B items. \freqlow & \ref{chap:recommendation} & Arch Design \\
17 & Daily training data volume for ads system (1B impressions/day). \freqlow & \ref{chap:recommendation} & Estimation \\
18 & High engagement but users churning. What is happening? \freqmed & \ref{chap:recommendation} & Debugging \\
19 & Efficient architecture for real-time video on edge devices. \freqlow & \ref{chap:efficient_architectures} & Arch Design \\
20 & Calculate training throughput (tokens/sec) for 7B on 8$\times$A100s. \freqlow & \ref{chap:training_optimization} & Estimation \\
21 & Design evaluation framework for LLM-based chatbot. \freqmed & \ref{chap:evaluation_metrics} & Arch Design \\
22 & Build vs buy vs open-source---ML infrastructure framework. \freqmed & \ref{chap:decision_frameworks} & Trade-off \\
23 & Estimate infrastructure cost: rec model at 100K QPS. \freqmed & \ref{chap:production_systems} & Estimation \\
24 & RLHF reward plateau while KL keeps climbing. What is happening? \freqlow & \ref{chap:reinforcement_learning} & Debugging \\
25 & Estimate compute cost of training CLIP on 400M pairs. \freqlow & \ref{chap:multimodal_learning} & Estimation \\
26 & LLM serving stack: 1K QPS, p99 $<$ 2s, 70B model. \freqhigh & \ref{chap:inference_optimization} & Arch Design \\
27 & Calculate throughput (tokens/sec) of 70B on 8$\times$H100s. \freqmed & \ref{chap:inference_optimization} & Estimation \\
28 & 500 concurrent users, 2s TTFT---size the infrastructure. \freqlow & \ref{chap:inference_optimization} & Estimation \\
29 & Prefill-optimized vs decode-optimized serving---disaggregate? \freqlow & \ref{chap:inference_optimization} & Trade-off \\
30 & Over-refusal: reduce without increasing safety risk. \freqmed & \ref{chap:safety_alignment} & Trade-off \\
31 & Steering vectors: control behavior without retraining. \freqlow & \ref{chap:safety_alignment} & Conceptual \\
32 & Design RAG for legal document search platform. \freqmed & \ref{chap:long_context_rag} & Arch Design \\
\end{longtable}

\begin{table}[H]
\centering
\caption{L7 question: RAG reasoning failure debugging}
\begin{tabular}{p{12cm}}
\toprule
33. Your RAG system works for factual queries but fails for reasoning queries. Diagnose. (\ref{chap:long_context_rag}, Debugging, \freqlow) \\
\bottomrule
\end{tabular}
\end{table}

%==============================================================================
\section{Questions by Type}
\label{sec:qi_by_type}
%==============================================================================

\subsection{Conceptual Questions}

52 questions testing understanding of \textit{what} and \textit{why} --- definitions, mechanisms, and design rationale.

\begin{longtable}{rp{8.5cm}ll}
\caption{Conceptual questions} \\
\toprule
\textbf{\#} & \textbf{Question} & \textbf{Ch} & \textbf{Diff} \\
\midrule
\endfirsthead
\toprule
\textbf{\#} & \textbf{Question} & \textbf{Ch} & \textbf{Diff} \\
\midrule
\endhead
\midrule
\multicolumn{4}{r}{\textit{Continued on next page}} \\
\endfoot
\bottomrule
\endlastfoot

1 & Explain SVD and three ML applications. & \ref{chap:mathematical_foundations} & L5 \\
2 & KL divergence: why not a true distance metric? & \ref{chap:mathematical_foundations} & L5 \\
3 & Eigenvalues and why they matter for training. & \ref{chap:mathematical_foundations} & L5 \\
4 & KL divergence asymmetry---when it matters in practice. & \ref{chap:mathematical_foundations} & L6 \\
5 & Double descent: test error decrease in overparameterized regime. & \ref{chap:learning_theory} & L5 \\
6 & Bias-variance tradeoff and deep learning. & \ref{chap:learning_theory} & L5 \\
7 & Emergent abilities: real or evaluation artifact? & \ref{chap:learning_theory} & L6 \\
8 & Grokking and what it tells us about generalization. & \ref{chap:learning_theory} & L6 \\
9 & Neural scaling laws: what scales and when they break. & \ref{chap:learning_theory} & L7 \\
10 & Softmax cross-entropy vs sigmoid BCE. & \ref{chap:loss_functions} & L5 \\
11 & SimCLR batch sizes and alternatives. & \ref{chap:loss_functions} & L5 \\
12 & Dying ReLU problem and modern solutions. & \ref{chap:activation_normalization} & L5 \\
13 & Flash Attention: faster with same computation? & \ref{chap:attention_transformers} & L5 \\
14 & RoPE position encoding: how and why preferred. & \ref{chap:attention_transformers} & L5 \\
15 & GPT-2 to LLaMA: what changed architecturally? & \ref{chap:attention_transformers} & L5 \\
16 & BPE and embedding quality interaction. & \ref{chap:embeddings} & L5 \\
17 & Embedding collapse: detect and prevent. & \ref{chap:embeddings} & L6 \\
18 & Cosine vs dot product vs L2 for embeddings. & \ref{chap:similarity_architectures} & L5 \\
19 & Augmentation policy for contrastive learning. & \ref{chap:self_supervised} & L5 \\
20 & BYOL without negatives: why no collapse? & \ref{chap:self_supervised} & L5 \\
21 & Decoder-only dominance despite encoder-decoder flexibility. & \ref{chap:nlp_architectures} & L6 \\
22 & RoPE vs learned vs sinusoidal positional embeddings. & \ref{chap:nlp_architectures} & L6 \\
23 & Effective receptive field of a CNN. & \ref{chap:vision_architectures} & L5 \\
24 & DETR and paradigm shift in object detection. & \ref{chap:vision_architectures} & L5 \\
25 & Classifier-free guidance in diffusion models. & \ref{chap:generative_models} & L5 \\
26 & Stable Diffusion end-to-end architecture. & \ref{chap:generative_models} & L5 \\
27 & Flow matching vs denoising diffusion. & \ref{chap:generative_models} & L7 \\
28 & Cold start problem in recommendations. & \ref{chap:recommendation} & L5 \\
29 & Recommendation model calibration: why needed? & \ref{chap:recommendation} & L5 \\
30 & ESMM and sample selection bias in CVR. & \ref{chap:recommendation} & L6 \\
31 & Heterogeneous graphs in GNNs. & \ref{chap:graph_neural_networks} & L5 \\
32 & Gradient accumulation: what and when. & \ref{chap:training_optimization} & L5 \\
33 & Negative transfer: detect and prevent. & \ref{chap:transfer_learning} & L5 \\
34 & Training-serving skew handling. & \ref{chap:production_systems} & L5 \\
35 & Concept drift vs data drift. & \ref{chap:production_systems} & L5 \\
36 & 95\% accuracy, unhappy stakeholders. & \ref{chap:evaluation_metrics} & L5 \\
37 & Evaluate a search ranking system. & \ref{chap:evaluation_metrics} & L5 \\
38 & Statistical significance of model improvements. & \ref{chap:evaluation_metrics} & L5 \\
39 & Accuracy as a metric: problems and when fine. & \ref{chap:evaluation_metrics} & L5 \\
40 & RLHF pipeline walkthrough. & \ref{chap:reinforcement_learning} & L5 \\
41 & Reward hacking: examples and mitigation. & \ref{chap:reinforcement_learning} & L5 \\
42 & CLIP vs SigLIP: what changed? & \ref{chap:multimodal_learning} & L5 \\
43 & Early vs late vs cross-attention fusion. & \ref{chap:multimodal_learning} & L5 \\
44 & Speculative decoding: when and when not. & \ref{chap:inference_optimization} & L5 \\
45 & Jailbreak attacks: how they work, defenses. & \ref{chap:safety_alignment} & L5 \\
46 & Hallucination detection and reduction. & \ref{chap:safety_alignment} & L5 \\
47 & Mechanistic interpretability: concrete circuit. & \ref{chap:safety_alignment} & L6 \\
48 & Steering vectors for behavior control. & \ref{chap:safety_alignment} & L7 \\
49 & RAG quality evaluation metrics. & \ref{chap:long_context_rag} & L5 \\
\end{longtable}

\subsection{Trade-off Questions}

64 questions testing the ability to \textit{compare} approaches and articulate when to choose one over another.

\begin{longtable}{rp{8.5cm}ll}
\caption{Trade-off questions} \\
\toprule
\textbf{\#} & \textbf{Question} & \textbf{Ch} & \textbf{Diff} \\
\midrule
\endfirsthead
\toprule
\textbf{\#} & \textbf{Question} & \textbf{Ch} & \textbf{Diff} \\
\midrule
\endhead
\midrule
\multicolumn{4}{r}{\textit{Continued on next page}} \\
\endfoot
\bottomrule
\endlastfoot

1 & Fixed compute budget: model size vs training data. & \ref{chap:learning_theory} & L6 \\
2 & 10B overparameterized model: reduce it? & \ref{chap:learning_theory} & L6 \\
3 & Train loss down, val plateau, 10$\times$ unlabeled data. & \ref{chap:learning_theory} & L6 \\
4 & 95\% accuracy but loss still decreasing---keep training? & \ref{chap:loss_functions} & L5 \\
5 & CE vs focal vs class-balanced for imbalanced data. & \ref{chap:loss_functions} & L5 \\
6 & Pointwise vs pairwise vs listwise ranking losses. & \ref{chap:loss_functions} & L6 \\
7 & GELU vs ReLU; SwiGLU vs GELU. & \ref{chap:activation_normalization} & L5 \\
8 & RMSNorm vs LayerNorm. & \ref{chap:activation_normalization} & L5 \\
9 & MHA vs GQA vs MQA. & \ref{chap:attention_transformers} & L5 \\
10 & Pre-norm vs Post-norm Transformers. & \ref{chap:attention_transformers} & L5 \\
11 & Transformer vs SSM (Mamba). & \ref{chap:attention_transformers} & L6 \\
12 & Flash Attention vs Linear Attention. & \ref{chap:attention_transformers} & L6 \\
13 & Static vs contextual embeddings. & \ref{chap:embeddings} & L5 \\
14 & Siamese vs Triplet vs Two-Tower. & \ref{chap:similarity_architectures} & L5 \\
15 & Cross-encoder vs bi-encoder. & \ref{chap:similarity_architectures} & L5 \\
16 & Hard negative mining strategies. & \ref{chap:similarity_architectures} & L6 \\
17 & SimCLR vs BYOL vs MAE. & \ref{chap:self_supervised} & L5 \\
18 & BERT vs GPT pre-training objectives. & \ref{chap:self_supervised} & L5 \\
19 & Encoder-only vs decoder-only vs encoder-decoder. & \ref{chap:nlp_architectures} & L5 \\
20 & BPE vs SentencePiece vs character-level. & \ref{chap:nlp_architectures} & L5 \\
21 & CNN vs ViT for 10K vs 10M images. & \ref{chap:vision_architectures} & L5 \\
22 & ResNet vs DenseNet vs Swin connections. & \ref{chap:vision_architectures} & L6 \\
23 & Diffusion vs GAN vs VAE vs autoregressive. & \ref{chap:generative_models} & L6 \\
24 & Two-Tower vs interaction-based for candidate gen. & \ref{chap:recommendation} & L5 \\
25 & Shared-bottom vs MMoE vs PLE. & \ref{chap:recommendation} & L6 \\
26 & GCN vs GraphSAGE vs GAT. & \ref{chap:graph_neural_networks} & L5 \\
27 & GNN vs simpler approaches. & \ref{chap:graph_neural_networks} & L6 \\
28 & GCN vs GraphSAGE vs GAT for production rec. & \ref{chap:graph_neural_networks} & L6 \\
29 & MoE vs dense scaling. & \ref{chap:efficient_architectures} & L6 \\
30 & Transformer vs Mamba for 100K tokens. & \ref{chap:efficient_architectures} & L6 \\
31 & INT8 vs INT4 vs FP8 quantization. & \ref{chap:efficient_architectures} & L6 \\
32 & Structured vs unstructured pruning. & \ref{chap:efficient_architectures} & L6 \\
33 & Adam vs SGD with momentum. & \ref{chap:training_optimization} & L5 \\
34 & Cosine decay vs constant LR with warmup. & \ref{chap:training_optimization} & L5 \\
35 & Data parallelism vs tensor vs pipeline parallelism. & \ref{chap:training_optimization} & L6 \\
36 & FP16 vs BF16 vs FP8 for training. & \ref{chap:training_optimization} & L6 \\
37 & LoRA vs QLoRA vs full FT vs prompt tuning. & \ref{chap:transfer_learning} & L6 \\
38 & Adapter vs LoRA vs prefix tuning. & \ref{chap:transfer_learning} & L5 \\
39 & LoRA vs full FT vs prompt tuning. & \ref{chap:transfer_learning} & L5 \\
40 & Shadow vs A/B test vs multi-armed bandit. & \ref{chap:production_systems} & L5 \\
41 & Precision 0.95, recall 0.30---what to do. & \ref{chap:evaluation_metrics} & L5 \\
42 & NDCG vs MAP vs MRR for ranking. & \ref{chap:evaluation_metrics} & L5 \\
43 & Latency 50ms to $<$10ms optimization tree. & \ref{chap:decision_frameworks} & L6 \\
44 & LLM fine-tune vs custom model (3 months, 3 engineers). & \ref{chap:decision_frameworks} & L6 \\
45 & GPU vs TPU vs CPU for inference. & \ref{chap:decision_frameworks} & L5 \\
46 & Build vs buy vs open-source for ML infra. & \ref{chap:decision_frameworks} & L7 \\
47 & BERT vs GPT for sentiment analysis. & \ref{chap:decision_frameworks} & L5 \\
48 & Gradient boosting vs neural nets for tabular. & \ref{chap:decision_frameworks} & L5 \\
49 & PPO vs DPO vs RLHF for alignment. & \ref{chap:reinforcement_learning} & L6 \\
50 & Constitutional AI vs RLHF vs DPO. & \ref{chap:reinforcement_learning} & L6 \\
51 & Process vs outcome reward models. & \ref{chap:reinforcement_learning} & L6 \\
52 & Adapter-based vs natively multimodal. & \ref{chap:multimodal_learning} & L6 \\
53 & Visual token reduction (576 tokens). & \ref{chap:multimodal_learning} & L6 \\
54 & GPTQ vs AWQ vs GGUF quantization. & \ref{chap:inference_optimization} & L6 \\
55 & Tensor vs pipeline parallelism for inference. & \ref{chap:inference_optimization} & L6 \\
56 & Prefill-optimized vs decode-optimized serving. & \ref{chap:inference_optimization} & L7 \\
57 & Over-refusal: reduce without increasing risk. & \ref{chap:safety_alignment} & L7 \\
58 & Chunking: fixed-size vs semantic vs document-aware. & \ref{chap:long_context_rag} & L5 \\
59 & Long context window vs RAG. & \ref{chap:long_context_rag} & L5 \\
60 & HyDE vs query rewriting vs decomposition. & \ref{chap:long_context_rag} & L6 \\
61 & Cross-encoder vs ColBERT vs bi-encoder reranking. & \ref{chap:long_context_rag} & L6 \\
\end{longtable}

\subsection{Estimation Questions}

28 questions testing back-of-envelope calculation, numerical fluency, and system sizing.

\begin{longtable}{rp{8.5cm}ll}
\caption{Estimation questions} \\
\toprule
\textbf{\#} & \textbf{Question} & \textbf{Ch} & \textbf{Diff} \\
\midrule
\endfirsthead
\toprule
\textbf{\#} & \textbf{Question} & \textbf{Ch} & \textbf{Diff} \\
\midrule
\endhead
\bottomrule
\endlastfoot

1 & 50K$\times$768 embedding memory (FP32/FP16/INT8). & \ref{chap:mathematical_foundations} & L5 \\
2 & FLOPs for (1024,4096)$\times$(4096,1024) matmul. & \ref{chap:mathematical_foundations} & L5 \\
3 & Chinchilla: compute-optimal model for $10^{22}$ FLOPs. & \ref{chap:learning_theory} & L6 \\
4 & 50K vocab, 768-dim embedding memory (FP16/INT8). & \ref{chap:embeddings} & L5 \\
5 & ANN index size and latency for 1B 768-dim embeddings. & \ref{chap:similarity_architectures} & L6 \\
6 & QPS for HNSW on 100M documents. & \ref{chap:similarity_architectures} & L6 \\
7 & QPS: HNSW 100M + cross-encoder top 100. & \ref{chap:similarity_architectures} & L6 \\
8 & BERT-base pre-training compute on 16GB text. & \ref{chap:self_supervised} & L6 \\
9 & 7B model inference latency: 100 tokens on A100. & \ref{chap:nlp_architectures} & L6 \\
10 & Memory to train 7B model (BS=1, seq=2048). & \ref{chap:attention_transformers} & L6 \\
11 & FLOPs for 7B Transformer forward pass. & \ref{chap:attention_transformers} & L6 \\
12 & KV cache memory: LLaMA 70B at 128K context. & \ref{chap:attention_transformers} & L6 \\
13 & Attention vs FFN FLOPs ratio at different seq lengths. & \ref{chap:attention_transformers} & L7 \\
14 & ResNet-50 vs ViT-B/16 FLOPs on 224$\times$224. & \ref{chap:vision_architectures} & L6 \\
15 & Stable Diffusion inference cost (512$\times$512, A100). & \ref{chap:generative_models} & L7 \\
16 & DLRM embedding memory (100M users, 10M items, 1K features). & \ref{chap:recommendation} & L6 \\
17 & Ads ranking: 100 candidates in $<$50ms constraints. & \ref{chap:recommendation} & L6 \\
18 & Daily training data volume (1B impressions/day). & \ref{chap:recommendation} & L7 \\
19 & 3-layer GCN memory on 10M-node graph. & \ref{chap:graph_neural_networks} & L6 \\
20 & MoE compute savings (8 experts, top-2 vs dense). & \ref{chap:efficient_architectures} & L6 \\
21 & GPU memory: 13B model LoRA vs full fine-tuning. & \ref{chap:training_optimization} & L6 \\
22 & Training throughput: 7B on 8$\times$A100s. & \ref{chap:training_optimization} & L7 \\
23 & LoRA trainable params: rank 16 on 7B model. & \ref{chap:transfer_learning} & L5 \\
24 & Infrastructure cost: rec model at 100K QPS. & \ref{chap:production_systems} & L7 \\
25 & A/B test sample size for 1\% CTR improvement. & \ref{chap:evaluation_metrics} & L6 \\
26 & RLHF memory (4 models). & \ref{chap:reinforcement_learning} & L6 \\
27 & CLIP training compute on 400M pairs. & \ref{chap:multimodal_learning} & L7 \\
28 & INT4 memory savings for 70B model. & \ref{chap:inference_optimization} & L6 \\
\end{longtable}

Plus 3 more L7 estimation questions: throughput of 70B on 8$\times$H100s (\ref{chap:inference_optimization}), 500 concurrent users infrastructure sizing (\ref{chap:inference_optimization}), and cost per query RAG vs long-context (\ref{chap:long_context_rag}).

\subsection{Debugging Questions}

38 questions testing systematic diagnosis of training, serving, and data pipeline failures.

\begin{longtable}{rp{8.5cm}ll}
\caption{Debugging questions} \\
\toprule
\textbf{\#} & \textbf{Question} & \textbf{Ch} & \textbf{Diff} \\
\midrule
\endfirsthead
\toprule
\textbf{\#} & \textbf{Question} & \textbf{Ch} & \textbf{Diff} \\
\midrule
\endhead
\midrule
\multicolumn{4}{r}{\textit{Continued on next page}} \\
\endfoot
\bottomrule
\endlastfoot

1 & Gradient norms 1000$\times$ larger in early layers. & \ref{chap:mathematical_foundations} & L6 \\
2 & Validation miscalibrated despite train loss decreasing. & \ref{chap:loss_functions} & L6 \\
3 & Contrastive loss jumps mid-training, doesn't recover. & \ref{chap:loss_functions} & L6 \\
4 & Training loss NaN after a few hundred steps (activation/norm). & \ref{chap:activation_normalization} & L5 \\
5 & Different behavior at training vs inference (normalization). & \ref{chap:activation_normalization} & L6 \\
6 & Validation accuracy lower than training (normalization). & \ref{chap:activation_normalization} & L5 \\
7 & Degrading perplexity at 32K despite 8K training. & \ref{chap:attention_transformers} & L6 \\
8 & All attention heads attending to same positions. & \ref{chap:attention_transformers} & L7 \\
9 & Embedding space: unrelated items cluster together. & \ref{chap:embeddings} & L6 \\
10 & Debug poor embedding quality. & \ref{chap:embeddings} & L5 \\
11 & Two-tower retrieves keyword matches, not relevant docs. & \ref{chap:similarity_architectures} & L5 \\
12 & Recall@10 good but precision@10 poor. & \ref{chap:similarity_architectures} & L6 \\
13 & Contrastive model: good for frequent, poor for rare items. & \ref{chap:similarity_architectures} & L7 \\
14 & SSL downstream plateaus despite 10$\times$ more data. & \ref{chap:self_supervised} & L6 \\
15 & Fine-tuned LLM: fluent but factually wrong. & \ref{chap:nlp_architectures} & L6 \\
16 & ViT 92\% ImageNet, 65\% on your domain. & \ref{chap:vision_architectures} & L6 \\
17 & Diffusion mode collapse for certain categories. & \ref{chap:generative_models} & L6 \\
18 & Rec model: great offline, poor online performance. & \ref{chap:recommendation} & L6 \\
19 & CTR AUC +0.5\% but revenue $-$3\%. & \ref{chap:recommendation} & L6 \\
20 & High engagement but users churning. & \ref{chap:recommendation} & L7 \\
21 & GNN performance degrades with more layers. & \ref{chap:graph_neural_networks} & L5 \\
22 & Quantized model fails on code/math domains. & \ref{chap:efficient_architectures} & L6 \\
23 & MoE routing collapse. & \ref{chap:efficient_architectures} & L6 \\
24 & Loss plateaus after initial descent. & \ref{chap:training_optimization} & L5 \\
25 & Training loss NaN after 1000 steps. & \ref{chap:training_optimization} & L5 \\
26 & Loss spikes every N steps, then recovers. & \ref{chap:training_optimization} & L6 \\
27 & Distributed training 60\% as fast with 8 GPUs. & \ref{chap:training_optimization} & L6 \\
28 & LoRA model catastrophically forgets general capabilities. & \ref{chap:transfer_learning} & L6 \\
29 & Good on test set, poor in production (after fine-tuning). & \ref{chap:transfer_learning} & L6 \\
30 & CTR dropped 5\% overnight. & \ref{chap:production_systems} & L6 \\
31 & Online metrics diverge after 2 weeks. & \ref{chap:production_systems} & L6 \\
32 & AUC 0.95 but ECE 0.15 (calibration). & \ref{chap:evaluation_metrics} & L6 \\
33 & Offline metrics up, online metrics down. & \ref{chap:evaluation_metrics} & L6 \\
34 & RLHF model becomes sycophantic. & \ref{chap:reinforcement_learning} & L6 \\
35 & RLHF: reward plateau, KL keeps climbing. & \ref{chap:reinforcement_learning} & L7 \\
36 & VLM hallucinates visual details not in image. & \ref{chap:multimodal_learning} & L6 \\
37 & KV cache memory bottleneck, dropping requests. & \ref{chap:inference_optimization} & L6 \\
38 & p99 latency 10$\times$ the p50. & \ref{chap:inference_optimization} & L6 \\
\end{longtable}

Plus 3 more: INT4 garbage on some prompts (\ref{chap:inference_optimization}, L6), speculative decoding slower than standard (\ref{chap:inference_optimization}, L6), RAG returns docs but LLM ignores them (\ref{chap:long_context_rag}, L6), and RAG fails for reasoning queries (\ref{chap:long_context_rag}, L7).

\subsection{Architecture Design Questions}

43 questions testing end-to-end system design, from problem formulation to production deployment.

\begin{longtable}{rp{8.5cm}ll}
\caption{Architecture Design questions} \\
\toprule
\textbf{\#} & \textbf{Question} & \textbf{Ch} & \textbf{Diff} \\
\midrule
\endfirsthead
\toprule
\textbf{\#} & \textbf{Question} & \textbf{Ch} & \textbf{Diff} \\
\midrule
\endhead
\midrule
\multicolumn{4}{r}{\textit{Continued on next page}} \\
\endfoot
\bottomrule
\endlastfoot

1 & Product recommendation with implicit feedback---loss. & \ref{chap:loss_functions} & L5 \\
2 & Embedding model for 100M doc semantic search---loss. & \ref{chap:loss_functions} & L5 \\
3 & Visual search for 10M products---loss function. & \ref{chap:loss_functions} & L6 \\
4 & Retrieval with click data---loss design. & \ref{chap:loss_functions} & L6 \\
5 & Multi-task model (classification + regression + ranking). & \ref{chap:loss_functions} & L7 \\
6 & Semantic search: choose and evaluate embeddings. & \ref{chap:embeddings} & L5 \\
7 & Embedding system for 100M product catalog. & \ref{chap:embeddings} & L6 \\
8 & Multi-stage retrieval pipeline design. & \ref{chap:similarity_architectures} & L6 \\
9 & Update embeddings without full ANN re-index. & \ref{chap:similarity_architectures} & L7 \\
10 & Design attention for 1M-token documents. & \ref{chap:attention_transformers} & L7 \\
11 & Modify Transformer for real-time streaming. & \ref{chap:attention_transformers} & L7 \\
12 & SSL pre-training pipeline for 10M unlabeled images. & \ref{chap:self_supervised} & L6 \\
13 & Medical imaging SSL pipeline (10M unlabeled, 10K labeled). & \ref{chap:self_supervised} & L6 \\
14 & Customer support classification (50 categories, multilingual). & \ref{chap:nlp_architectures} & L6 \\
15 & Multilingual doc understanding (50+ languages). & \ref{chap:nlp_architectures} & L7 \\
16 & Edge device vision backbone ($<$10W). & \ref{chap:vision_architectures} & L6 \\
17 & Visual search: 100M images, sub-200ms. & \ref{chap:vision_architectures} & L6 \\
18 & Product photography image generation pipeline. & \ref{chap:generative_models} & L6 \\
19 & Video streaming rec system (100M users, 1M videos). & \ref{chap:recommendation} & L6 \\
20 & E-commerce product ads ranking model. & \ref{chap:recommendation} & L6 \\
21 & Real-time personalized notification system. & \ref{chap:recommendation} & L7 \\
22 & Embedding infrastructure for 1B items. & \ref{chap:recommendation} & L7 \\
23 & GNN-based fraud detection system. & \ref{chap:graph_neural_networks} & L6 \\
24 & On-device ML ($<$100ms, $<$50MB). & \ref{chap:efficient_architectures} & L6 \\
25 & KD pipeline: 70B to 7B. & \ref{chap:efficient_architectures} & L6 \\
26 & Real-time video on edge devices. & \ref{chap:efficient_architectures} & L7 \\
27 & Scale 70B model training across 256 GPUs. & \ref{chap:training_optimization} & L6 \\
28 & Fine-tune 70B on single A100 80GB. & \ref{chap:transfer_learning} & L5 \\
29 & Fine-tuning pipeline for 10K-example domain adaptation. & \ref{chap:transfer_learning} & L6 \\
30 & Search ranking infra at 10K QPS. & \ref{chap:production_systems} & L6 \\
31 & A/B test setup for recommendation model. & \ref{chap:production_systems} & L6 \\
32 & ML monitoring system design. & \ref{chap:production_systems} & L6 \\
33 & Feature store: batch vs streaming architecture. & \ref{chap:production_systems} & L6 \\
34 & LLM chatbot evaluation framework. & \ref{chap:evaluation_metrics} & L7 \\
35 & Content recommendation for news app. & \ref{chap:decision_frameworks} & L6 \\
36 & Content moderation system. & \ref{chap:decision_frameworks} & L6 \\
37 & Reward model training pipeline. & \ref{chap:reinforcement_learning} & L6 \\
38 & Multimodal content understanding for social media. & \ref{chap:multimodal_learning} & L6 \\
39 & Multimodal product search. & \ref{chap:multimodal_learning} & L6 \\
40 & Add audio to existing VLM. & \ref{chap:multimodal_learning} & L6 \\
41 & LLM serving stack: 1K QPS, p99 $<$ 2s, 70B. & \ref{chap:inference_optimization} & L7 \\
42 & Content filtering pipeline ($<$50ms). & \ref{chap:safety_alignment} & L6 \\
43 & Hallucination detection system. & \ref{chap:safety_alignment} & L6 \\
\end{longtable}

Plus: red teaming program (\ref{chap:safety_alignment}, L6), legal RAG system (\ref{chap:long_context_rag}, L7), enterprise 10M-doc RAG (\ref{chap:long_context_rag}, L6).

\subsection{First Principles Questions}

23 questions testing the ability to \textit{derive} results, explain \textit{why} something works from mathematical or theoretical foundations.

\begin{longtable}{rp{8.5cm}ll}
\caption{First Principles questions} \\
\toprule
\textbf{\#} & \textbf{Question} & \textbf{Ch} & \textbf{Diff} \\
\midrule
\endfirsthead
\toprule
\textbf{\#} & \textbf{Question} & \textbf{Ch} & \textbf{Diff} \\
\midrule
\endhead
\bottomrule
\endlastfoot

1 & Chain rule and backprop: walk through 3 layers. & \ref{chap:mathematical_foundations} & L5 \\
2 & Derive softmax CE gradient w.r.t.\ logits. & \ref{chap:mathematical_foundations} & L6 \\
3 & Rank of attention matrix and efficiency. & \ref{chap:mathematical_foundations} & L7 \\
4 & Why LLMs generalize with more params than examples. & \ref{chap:learning_theory} & L7 \\
5 & Why Transformers don't use BatchNorm (first principles). & \ref{chap:activation_normalization} & L6 \\
6 & Derive BatchNorm incompatibility with Transformers. & \ref{chap:activation_normalization} & L6 \\
7 & Why attention scales by $1/\sqrt{d_k}$. & \ref{chap:attention_transformers} & L5 \\
8 & Self-attention permutation equivariance. & \ref{chap:attention_transformers} & L5 \\
9 & Next-token vs MLM: why autoregressive won. & \ref{chap:self_supervised} & L6 \\
10 & InfoNCE loss and mutual information connection. & \ref{chap:self_supervised} & L6 \\
11 & In-context learning mechanistically. & \ref{chap:nlp_architectures} & L7 \\
12 & DPO loss: eliminates separate reward model. & \ref{chap:nlp_architectures} & L6 \\
13 & Reparameterization trick for VAEs. & \ref{chap:generative_models} & L5 \\
14 & ELBO derivation and blurry VAE images. & \ref{chap:generative_models} & L5 \\
15 & WGAN loss and Wasserstein distance. & \ref{chap:generative_models} & L6 \\
16 & WL expressivity of message-passing GNNs. & \ref{chap:graph_neural_networks} & L6 \\
17 & Why LoRA works: low-rank hypothesis. & \ref{chap:transfer_learning} & L6 \\
18 & KL penalty in RLHF: why it matters. & \ref{chap:reinforcement_learning} & L5 \\
19 & CLIP contrastive loss: shared embedding space. & \ref{chap:multimodal_learning} & L5 \\
\end{longtable}

%==============================================================================
\section{Questions by Topic}
\label{sec:qi_by_topic}
%==============================================================================

\begin{table}[H]
\centering
\caption{Question counts by broad topic area}
\begin{tabularx}{\textwidth}{Xlr}
\toprule
\textbf{Topic Area} & \textbf{Chapters} & \textbf{Questions} \\
\midrule
Math \& Theory & Ch~\ref{chap:mathematical_foundations}--\ref{chap:learning_theory} & 20 \\
Loss Functions & Ch~\ref{chap:loss_functions} & 12 \\
Activations \& Normalization & Ch~\ref{chap:activation_normalization} & 8 \\
Attention \& Transformers & Ch~\ref{chap:attention_transformers} & 18 \\
Embeddings \& Retrieval & Ch~\ref{chap:embeddings}--\ref{chap:similarity_architectures} & 20 \\
Self-Supervised Learning & Ch~\ref{chap:self_supervised} & 10 \\
NLP Architectures & Ch~\ref{chap:nlp_architectures} & 10 \\
Vision Architectures & Ch~\ref{chap:vision_architectures} & 8 \\
Generative Models & Ch~\ref{chap:generative_models} & 10 \\
Recommendation Systems & Ch~\ref{chap:recommendation} & 15 \\
Graph Neural Networks & Ch~\ref{chap:graph_neural_networks} & 8 \\
Efficient Architectures & Ch~\ref{chap:efficient_architectures} & 10 \\
Training Optimization & Ch~\ref{chap:training_optimization} & 12 \\
Transfer Learning & Ch~\ref{chap:transfer_learning} & 10 \\
Production Systems & Ch~\ref{chap:production_systems} & 10 \\
Evaluation Metrics & Ch~\ref{chap:evaluation_metrics} & 10 \\
Decision Frameworks & Ch~\ref{chap:decision_frameworks} & 8 \\
Reinforcement Learning \& Alignment & Ch~\ref{chap:reinforcement_learning} & 10 \\
Multimodal Learning & Ch~\ref{chap:multimodal_learning} & 10 \\
Inference Optimization & Ch~\ref{chap:inference_optimization} & 12 \\
Safety \& Alignment & Ch~\ref{chap:safety_alignment} & 8 \\
Long Context \& RAG & Ch~\ref{chap:long_context_rag} & 10 \\
\midrule
\textbf{Total} & & \textbf{249} \\
\bottomrule
\end{tabularx}
\end{table}

\subsection{LLM Fundamentals}
\textit{The most heavily tested topic area in 2024--2025 interviews.}

\begin{table}[H]
\centering
\caption{LLM-focused questions across chapters}
\begin{tabularx}{\textwidth}{Xllr}
\toprule
\textbf{Question} & \textbf{Ch} & \textbf{Diff} & \textbf{Freq} \\
\midrule
Flash Attention: why faster? & \ref{chap:attention_transformers} & L5 & \freqhigh \\
MHA vs GQA vs MQA & \ref{chap:attention_transformers} & L5 & \freqhigh \\
GPT-2 to LLaMA changes & \ref{chap:attention_transformers} & L5 & \freqhigh \\
KV cache memory for 70B at 128K & \ref{chap:attention_transformers} & L6 & \freqhigh \\
Encoder vs decoder vs enc-dec & \ref{chap:nlp_architectures} & L5 & \freqhigh \\
Decoder-only dominance & \ref{chap:nlp_architectures} & L6 & \freqhigh \\
RLHF pipeline walkthrough & \ref{chap:reinforcement_learning} & L5 & \freqhigh \\
PPO vs DPO vs RLHF & \ref{chap:reinforcement_learning} & L6 & \freqhigh \\
DPO loss derivation & \ref{chap:nlp_architectures} & L6 & \freqhigh \\
Speculative decoding & \ref{chap:inference_optimization} & L5 & \freqhigh \\
KV cache bottleneck solutions & \ref{chap:inference_optimization} & L6 & \freqhigh \\
LLM serving stack (1K QPS, 70B) & \ref{chap:inference_optimization} & L7 & \freqhigh \\
Jailbreak attacks and defenses & \ref{chap:safety_alignment} & L5 & \freqhigh \\
Hallucination detection/reduction & \ref{chap:safety_alignment} & L5 & \freqhigh \\
RAG vs long context & \ref{chap:long_context_rag} & L5 & \freqhigh \\
\bottomrule
\end{tabularx}
\end{table}

\subsection{Ads, Ranking, and Recommendations}
\textit{Critical for Meta, Google, Amazon, and TikTok interviews.}

\begin{table}[H]
\centering
\caption{Ads/Ranking/RecSys questions}
\begin{tabularx}{\textwidth}{Xllr}
\toprule
\textbf{Question} & \textbf{Ch} & \textbf{Diff} & \textbf{Freq} \\
\midrule
Video streaming rec system (100M users) & \ref{chap:recommendation} & L6 & \freqhigh \\
E-commerce ads ranking model & \ref{chap:recommendation} & L6 & \freqhigh \\
Cold start problem & \ref{chap:recommendation} & L5 & \freqhigh \\
Two-Tower vs interaction-based & \ref{chap:recommendation} & L5 & \freqhigh \\
DLRM embedding memory estimation & \ref{chap:recommendation} & L6 & \freqhigh \\
CTR AUC up but revenue down & \ref{chap:recommendation} & L6 & \freqhigh \\
Recommendation calibration & \ref{chap:recommendation} & L5 & \freqhigh \\
Multi-stage retrieval pipeline & \ref{chap:similarity_architectures} & L6 & \freqhigh \\
Pointwise vs pairwise vs listwise losses & \ref{chap:loss_functions} & L6 & \freqhigh \\
Offline up, online down (5 reasons) & \ref{chap:evaluation_metrics} & L6 & \freqhigh \\
\bottomrule
\end{tabularx}
\end{table}

\subsection{Training and Optimization}
\textit{Essential for ML infrastructure and training-focused roles.}

\begin{table}[H]
\centering
\caption{Training and optimization questions}
\begin{tabularx}{\textwidth}{Xllr}
\toprule
\textbf{Question} & \textbf{Ch} & \textbf{Diff} & \textbf{Freq} \\
\midrule
Adam vs SGD with momentum & \ref{chap:training_optimization} & L5 & \freqhigh \\
Loss plateaus debugging & \ref{chap:training_optimization} & L5 & \freqhigh \\
Training loss NaN debugging & \ref{chap:training_optimization} & L5 & \freqhigh \\
Scale 70B across 256 GPUs & \ref{chap:training_optimization} & L6 & \freqhigh \\
Data/tensor/pipeline parallelism & \ref{chap:training_optimization} & L6 & \freqhigh \\
13B GPU memory: LoRA vs full FT & \ref{chap:training_optimization} & L6 & \freqhigh \\
LoRA vs QLoRA vs full FT vs prompt tuning & \ref{chap:transfer_learning} & L6 & \freqhigh \\
FP16 vs BF16 vs FP8 & \ref{chap:training_optimization} & L6 & \freqmed \\
Chinchilla compute-optimal scaling & \ref{chap:learning_theory} & L6 & \freqhigh \\
\bottomrule
\end{tabularx}
\end{table}

\subsection{Production and Serving}
\textit{Differentiates ML engineers from ML researchers.}

\begin{table}[H]
\centering
\caption{Production and serving questions}
\begin{tabularx}{\textwidth}{Xllr}
\toprule
\textbf{Question} & \textbf{Ch} & \textbf{Diff} & \textbf{Freq} \\
\midrule
Training-serving skew & \ref{chap:production_systems} & L5 & \freqhigh \\
CTR dropped 5\% overnight & \ref{chap:production_systems} & L6 & \freqhigh \\
A/B test setup for rec model & \ref{chap:production_systems} & L6 & \freqhigh \\
Online diverges from offline after 2 weeks & \ref{chap:production_systems} & L6 & \freqhigh \\
Shadow vs A/B vs bandit & \ref{chap:production_systems} & L5 & \freqhigh \\
Search ranking infra at 10K QPS & \ref{chap:production_systems} & L6 & \freqhigh \\
Tensor vs pipeline parallelism for inference & \ref{chap:inference_optimization} & L6 & \freqhigh \\
INT4 quantization savings for 70B & \ref{chap:inference_optimization} & L6 & \freqhigh \\
Content filtering $<$50ms & \ref{chap:safety_alignment} & L6 & \freqhigh \\
\bottomrule
\end{tabularx}
\end{table}
